\section*{الفصل \ref{ch:b2:reaction-enthalpy} --- أنتالبيات التفاعل}

\begin{solution}{exo:b2:reaction-enthalpy:1}
$\Delta_r H^\circ = 3(-393.51) + 4(-285.83) - (-104.7) - 5(0) =
\qty{-2219.2}{kJ/mol}$، وهو سالب: التفاعل \omterm{def:b2:reaction-enthalpy:exothermic}{ناشر للحرارة}. (وأنتالبي الاحتراق
المقيس هو نفسه، \qty{-2219.2}{kJ/mol}.)
\end{solution}

\begin{solution}{exo:b2:reaction-enthalpy:2}
الميثان: $890.3/16.04 = \qty{55.5}{kJ/g}$، أي \qty{55.5}{MJ} لكل كيلوغرام.
ثنائي الهيدروجين: $285.83/2.016 = \qty{141.8}{kJ/g}$،
أي \qty{141.8}{MJ/kg}: أكثر بنحو 2.6 مرة لكل كيلوغرام (لكنه أقل بكثير لكل
لتر، لأن الغاز خفيف).
\end{solution}

\begin{solution}{exo:b2:reaction-enthalpy:3}
$\Delta\nu_{\text{غاز}} = 6 - 7.5 = -1.5$، ومنه $\Delta_r H = \Delta_r U +
\Delta\nu_{\text{غاز}}RT = -3264 - 1.5 \times 8.314 \times 298.15 \times
10^{-3} = \qty{-3267.7}{kJ/mol}$، وهو يتفق مع القيمة المحسوبة من أنتالبيات
التكوين، \qty{-3267.6}{kJ/mol}.
\end{solution}

\begin{solution}{exo:b2:reaction-enthalpy:4}
التفاعل المستهدف هو التفاعل الأول ناقص ضعف الثاني:
$\Delta_r H^\circ = -393.5 - 2(-283.0) = \qty{+172.5}{kJ/mol}$. فهو \omterm{def:b2:reaction-enthalpy:exothermic}{ماص
للحرارة}: يبرد فحم الكوك الساخن وهو يحوّل ثنائي أكسيد الكربون إلى أحادي
أكسيد الكربون.
\end{solution}

\begin{solution}{exo:b2:reaction-enthalpy:5}
الروابط المكسورة: C--H وCl--Cl، $416 + 243 = \qty{659}{kJ/mol}$؛ والروابط
المتكوّنة: C--Cl وH--Cl، $350 + 432 = \qty{782}{kJ/mol}$: التقدير
\qty{-123}{kJ/mol}. ومن أنتالبيات التكوين: $-83.7 - 92.3 + 74.9 =
\qty{-101.1}{kJ/mol}$. ويأتي الفرق البالغ \qty{22}{kJ/mol} من قيمة C--Cl
المتوسطة، وهي معدّل على عدة جزيئات، ومن رابطة C--H المكسورة في الميثان،
وهي أقوى من متوسط الروابط الأربع.
\end{solution}

\begin{solution}{exo:b2:reaction-enthalpy:6}
عند \qty{298}{K}: $-241.83 + 285.83 = \qty{44.00}{kJ/mol}$. وبقانون كيرشوف مع
$\Delta C_p = 33.59 - 75.35 = \qty{-41.76}{J/(K.mol)}$:
$44.00 - 0.04176 \times 101.85 = \qty{39.75}{kJ/mol}$ عند \qty{400}{K}. وتعطي
الجداول $-242.85 + 282.59 = \qty{39.74}{kJ/mol}$: فالسعات الحرارية الثابتة
ممتازة على هذه الفترة.
\end{solution}

\begin{solution}{exo:b2:reaction-enthalpy:7}
التأين: $4.341 \times 96.485 = \qty{418.8}{kJ/mol}$؛ واكتساب الإلكترون:
$-3.613 \times 96.485 = \qty{-348.6}{kJ/mol}$. \omterm{def:b2:reaction-enthalpy:lattice-enthalpy}{أنتالبي الشبكة} $= 436.7 +
89.0 + 121.3 + 418.8 - 348.6 = \qty{717}{kJ/mol}$، وهو أقل من
\qty{787}{kJ/mol} لكلوريد الصوديوم: فالأيون \ce{K+} أكبر من \ce{Na+}،
فيتباعد الأيونان أكثر ويتجاذبان أقل.
\end{solution}

\begin{solution}{exo:b2:reaction-enthalpy:8}
السعة الحرارية للمسعر: $C = 2.00/0.418 = \qty{4.785}{kJ/K}$.
الحرارة المحررة: $4.785 \times 0.583 = \qty{2.789}{kJ}$، من أجل $\xi =
\qty{0.0500}{mol}$ (الحمض هو المحدِّد، والقاعدة بإفراط طفيف):
$\Delta_r H = -2.789/0.0500 = \qty{-55.8}{kJ/mol}$.
\end{solution}

\begin{solution}{exo:b2:reaction-enthalpy:9}
\ce{C8H18(l) + 25/2O2(g) -> 8CO2(g) + 9H2O(l)}: $\Delta_r H^\circ =
8(-393.51) + 9(-285.83) + 250.3 = \qty{-5470.3}{kJ/mol}$؛ ومع $M =
\qty{114.2}{g/mol}$ نجد \qty{47.9}{MJ/kg}. ثنائي أكسيد الكربون: $8 \times 44.01 =
\qty{352.1}{g}$ مقابل \qty{5.470}{MJ}، أي \qty{64.4}{g/MJ}، مقابل
$44.01/0.8903 = \qty{49.4}{g/MJ}$ للميثان: ففي الميثان هيدروجين أكثر لكل
ذرة كربون، واحتراق الهيدروجين يحرر حرارة دون ثنائي أكسيد الكربون.
\end{solution}

\begin{solution}{exo:b2:reaction-enthalpy:10}
نواتج \ce{CH4 + 2O2 -> CO2 + 2H2O(g)}: $C_p = 37.13 + 2 \times 33.59 =
\qty{104.3}{J/K}$؛ $q = \qty{802.3}{kJ}$؛ $T_f = 298 + 802\,300/104.3 \approx
\qty{7990}{K}$، أي نحو ثلاثة أضعاف القيمة في الهواء، لأنه لا يوجد نيتروجين
يأخذ نصيبه من الحرارة. واللهب الحقيقي في الأكسجين أبرد بكثير: فقبل بلوغ
\qty{7990}{K} بزمن طويل، يتفكك ثنائي أكسيد الكربون والماء إلى أحادي أكسيد
الكربون وهيدروجين وأكسجين وذرات، في تفاعلات ماصة للحرارة تمتص معظم
الحرارة؛ كما أن السعات الحرارية تزداد مع
$T$.
\end{solution}

\begin{solution}{exo:b2:reaction-enthalpy:11}
$\Delta_r H^\circ(700) = \Delta_r H^\circ(298) + \int_{298}^{700}(\Delta_r a +
\Delta_r b\,T)\,\dd T = -91.88 + [-60.5 \times 401.85 + 0.0270 \times
(700^2 - 298.15^2)] \times 10^{-3} = -91.88 - 24.31 + 10.83 =
\qty{-105.4}{kJ/mol}$، أي في حدود \qty{0.2}{kJ/mol} من قيمة الجداول
(\qty{-105.2}{kJ/mol})، في حين أخطأ التقدير بسعة $C_p$ ثابتة بمقدار
\qty{4.5}{kJ/mol}.
\end{solution}

\begin{solution}{exo:b2:reaction-enthalpy:12}
$\Delta_r H^\circ = 2 \times 90.29 = \qty{+180.6}{kJ/mol}$ (\omterm{def:b2:reaction-enthalpy:exothermic}{ماص للحرارة}).
$\Delta_r C_p^\circ = 2(29.85) - 29.12 - 29.38 = \qty{1.2}{J/(K.mol)}$، ومنه
فإن \omterm{def:b2:reaction-enthalpy:reaction-quantity}{أنتالبي التفاعل} يتغير بين \qty{298}{K} و\qty{2000}{K} بنحو
$1.2 \times 1702 = \qty{2}{kJ/mol}$، أي واحد في المئة. والتفاعل \omterm{def:b2:reaction-enthalpy:exothermic}{الماص للحرارة}
تحابيه درجة الحرارة المرتفعة (\cref{ch:b2:equilibrium-shifts}):
فالهواء البارد لا يكاد يصنع أحادي أكسيد النيتروجين، وأسخن اللهب يصنع
أكثره، ولهذا تُعدّ درجة حرارة الاحتراق رافعة للحد من هذا
الملوِّث.
\end{solution}

\begin{solution}{pb:b2:reaction-enthalpy:1}
\textbf{1.} \ce{C4H10(g) + 13/2O2(g) -> 4CO2(g) + 5H2O(l)}.
\textbf{2.} $\Delta_r H^\circ = 4(-393.51) + 5(-285.83) + 125.6 =
\qty{-2877.6}{kJ/mol}$.
\textbf{3.} مع بخار الماء، $4(-393.51) + 5(-241.83) + 125.6 =
\qty{-2657.6}{kJ/mol}$؛ والفرق، \qty{220.0}{kJ}، هو أنتالبي تبخير خمسة
مولات من الماء (\qty{44.0}{kJ/mol} لكل منها).
\textbf{4.} $2877.6/58.12 = \qty{49.5}{kJ/g}$ و$2657.6/58.12 =
\qty{45.7}{kJ/g}$.
\textbf{5.} الثانية: يغادر الماء بخارًا وتضيع حرارة تكثفه.
\textbf{6.} $230 \times 45.7 = \qty{1.05e4}{kJ}$، أي نحو \qty{10.5}{MJ}.
\textbf{7.} عند حجم ثابت تكون الحرارة المستقبَلة $\Delta U$ (لا شغل
للضغط): فالقنبلة تقيس $\Delta_r U$.
\textbf{8.} $\Delta\nu_{\text{غاز}} = 4 - 1 - 6.5 = -3.5$ (الماء
سائل).
\textbf{9.} $\Delta_r U = \Delta_r H - \Delta\nu_{\text{غاز}}RT = -2877.6 +
3.5 \times 2.479 = \qty{-2868.9}{kJ/mol}$.
\textbf{10.} $\xi = 0.500/58.12 = \qty{8.60e-3}{mol}$؛ والحرارة $8.60 \times
10^{-3} \times 2868.9 = \qty{24.68}{kJ}$؛ $\Delta T = 24.68/10.00 =
\qty{2.468}{K}$.
\textbf{11.} الحرارة $10.00 \times 2.461 = \qty{24.61}{kJ}$: $\Delta_r U =
-24.61/(8.603 \times 10^{-3}) = \qty{-2860.6}{kJ/mol}$ و$\Delta_r H =
-2860.6 - 8.7 = \qty{-2869.3}{kJ/mol}$، أي أدنى من قيمة الجداول بمقدار
\qty{0.29}{\percent} بالقيمة المطلقة: حرارة فقدها المسعر أو احتراق غير تام
بعض الشيء.
\textbf{12.} $n = 1000/18.02 = \qty{55.5}{mol}$؛ $Q = 55.5 \times 75.35
\times 80 = \qty{334.5}{kJ}$.
\textbf{13.} الحرارة المحررة $16.0 \times 45.72 = \qty{731.5}{kJ}$؛ والكفاءة
$334.5/731.5 = 0.46$.
\textbf{14.} إلى غازات الاحتراق الساخنة التي تفلت حول القدر، وإلى الهواء
الذي يسخّنه اللهب، وإلى جدران القدر، وإلى الإشعاع.
\textbf{15.} $230 \times 45.72 \times 0.457/334.5 = \qty{14.4}{L}$.
\textbf{16.} يرافق \qty{6.5}{mol} من \ce{O2} مقدار $6.5 \times 78.08/20.95 =
\qty{24.23}{mol}$ من \ce{N2} و$6.5 \times 0.93/20.95 = \qty{0.289}{mol}$
من \ce{Ar}.
\textbf{17.} \qty{4}{mol} من \ce{CO2}، و\qty{5}{mol} من \ce{H2O(g)}،
و\qty{24.23}{mol} من \ce{N2}، و\qty{0.289}{mol} من \ce{Ar}.
\textbf{18.} اللهب كظوم ومتساوي الضغط، فيكون $\Delta H = 0$. وعلى طول
مسار مكوَّن من التفاعل عند \qty{298}{K} ($\Delta H_1 =
\qty{-2657.6}{kJ}$) متبوعًا بتسخين هذه النواتج، يكون $\Delta H_2 =
-\Delta H_1 = \qty{+2657.6}{kJ}$: فالحرارة كلها تذهب إلى الغازات.
\textbf{19.} $C_p = 4(37.13) + 5(33.59) + 24.23(29.12) + 0.289(20.79) =
\qty{1028}{J/K}$.
\textbf{20.} $T_f = 298 + 2\,657\,600/1028 = \qty{2883}{K}$.
\textbf{21.} $T_f = 2300 + 100 \times (2657.6 - 2515.7)/(2655.7 - 2515.7) =
\qty{2401}{K}$.
\textbf{22.} تزداد السعات الحرارية للغازات مع درجة الحرارة (بنحو الربع
بين \qty{298}{K} و\qty{2400}{K}): فقيم درجة حرارة الغرفة تجعل تسخين
الغازات أرخص مما هو في الواقع.
\textbf{23.} فوق \qty{2000}{K} يتفكك ثنائي أكسيد الكربون والماء جزئيًا
(تفكك \omterm{def:b2:reaction-enthalpy:exothermic}{ماص للحرارة})، ويشع اللهب، ولا يكون الامتزاج بالهواء تامًا أبدًا.
\textbf{24.} \omterm{def:b2:reaction-enthalpy:flame-temperature}{درجة حرارة اللهب الكظومة} للبوتان في الهواء هي
$\boldsymbol{T_{\text{ad}} \approx \qty{2.40e3}{K}}$.
\end{solution}
